\chapter{Introduction}
\label{chap:introduction}

\section{Context: why lunar subsurface thermal models still matter}

Sixty years after the first Surveyor landers, the Moon's thermal
environment remains the single most important \emph{boundary condition}
for every volatile-stability, in-situ-resource-utilisation (ISRU), and
permanently shadowed region (PSR) study in the run-up to the Artemis
programme.  The Diviner radiometer on the Lunar Reconnaissance Orbiter
(LRO) has mapped top-of-regolith bolometric temperatures at
\SI{237}{\meter} pixel scale, but the \emph{subsurface} is
intrinsically opaque: we see it only through the handful of in-situ
probes emplaced by Apollo~15/17, Chang'E-4, and Chandrayaan-3's ChaSTE
instrument.

Every published lunar thermal model must therefore survive the same
four sanity checks:
\begin{enumerate}
  \item does it reproduce the Apollo HFE deep-sensor equilibrium
        temperatures ($\sim$\SI{250}{\kelvin}) at each site?
  \item does it recover the diurnal surface-temperature amplitude
        observed by Diviner and Chang'E-4?
  \item can it be extended to a polar geometry where self-shadowing and
        ice-coupled conductivity dominate?
  \item does it conserve energy down to a realistic geothermal flux at
        the bottom boundary?
\end{enumerate}
\noindent The \texttt{Lunar-V2} pipeline was written to answer all four
simultaneously and to do so in code that a graduate student --- without
legacy Fortran infrastructure --- can run on a laptop in a handful of
minutes per site.

\section{Scope of this manual}

\Cref{fig:pipeline-flow} summarises the four mission phases covered in
this book.  Phase~0 (\emph{quickstart}) is not a research phase; it is
a smoke-test notebook that verifies the \code{lunar} package imports
cleanly on a freshly cloned repository.  The chapters that follow treat
Phases~1 through~4 in order, each ending with a bridge to the next.

\begin{figure}[H]
\centering
\begin{tikzpicture}[node distance=1.0cm and 1.6cm]
  \node[doneBlock]   (p0) {\textbf{Phase 0}\\Quickstart\\\tiny shipped};
  \node[doneBlock,   right=of p0]   (p1) {\textbf{Phase 1}\\Point-source\\validation\\\tiny shipped};
  \node[activeBlock, right=of p1]   (p2) {\textbf{Phase 2}\\Illumination\\coupling\\\tiny in progress};
  \node[plannedBlock,right=of p2]   (p3) {\textbf{Phase 3}\\Polar maps\\+ ice\\\tiny planned};
  \node[plannedBlock,right=of p3]   (p4) {\textbf{Phase 4}\\Manuscript\\+ thesis\\\tiny planned};
  \draw[arrowThick] (p0) -- (p1);
  \draw[arrowThick] (p1) -- (p2);
  \draw[arrowThick] (p2) -- (p3);
  \draw[arrowThick] (p3) -- (p4);
\end{tikzpicture}
\caption[Pipeline phase flow]{The five-stage \texttt{Lunar-V2} / TSUKIMI
  pipeline.  Green = shipped; amber = in progress; grey = planned.
  The book covers Phases~1--4.}
\label{fig:pipeline-flow}
\end{figure}

\begin{concept}
The pipeline is deliberately \emph{linear}.  Each phase receives the
previous phase's validated output (parameter set, DEM mosaic, solver
configuration) and returns a new validated deliverable.  No phase
short-circuits another; failure to meet a phase's acceptance criterion
blocks the next.
\end{concept}

\section{Code base at a glance}

The repository\footnote{\repo, branch \branch.} follows a conventional
Python scientific layout:

\begin{lstlisting}[caption={Top-level directory layout (abridged).},label={lst:repo-layout},language=bash]
Lunar-V2/
|-- lunar/                        # importable package
|   |-- solver_cn.py              # 1-D Crank-Nicolson core
|   |-- properties.py             # regolith property closures
|   |-- insolation.py             # solar forcing + SPICE wrapper
|   |-- apollo_helpers.py         # Phase-1 validation helpers
|   `-- pixel_column.py           # per-pixel column driver
|-- notebooks/
|   |-- phase0_quickstart/
|   |-- phase1_validation/        # 01_apollo_validation.ipynb + ...
|   |-- phase2_illumination/      # DEM + horizon + shadow
|   |-- phase3_polar_maps/        # Diviner + ice stability
|   `-- phase4_manuscript/        # PSJ figure bundle
|-- scripts/                      # standalone CLI utilities
|-- output/
|   |-- figures/                  # publication-ready PDF + PNG
|   |-- gifs/                     # thermal-evolution animations
|   `-- docs/                     # compiled PDF roadmap
`-- docs/roadmap/                 # this book (LaTeX source)
\end{lstlisting}

Every notebook is self-installing and self-fetching: \code{pip}
dependencies, SPICE kernels, and PDS HFE archives are downloaded on
first \emph{Run All}.  This is a deliberate design choice that
supports the reproducibility requirement (\cref{sec:novelty-reproducibility}).

\section{Nomenclature}

Throughout the book, we adopt SI units exclusively.  The symbol
conventions used in the equations are summarised in
\cref{tab:symbols}.

\begin{table}[H]
\centering
\small
\begin{tabularx}{\linewidth}{l l Y}
\toprule
Symbol & Unit & Meaning \\
\midrule
$T$        & K         & Temperature \\
$\rho$     & \si{\kilogram\per\meter\cubed} & Bulk density \\
$c_p$      & \si{\joule\per\kilogram\per\kelvin} & Specific heat capacity \\
$k$        & \si{\watt\per\meter\per\kelvin} & Thermal conductivity \\
$\Ksurf,\Kdeep$ & \si{\watt\per\meter\per\kelvin} & Surface / deep $k$ endpoints \\
$\Hscale$  & \si{\meter} & Hayne e-folding scale for $k$ \\
$\Qb$      & \si{\milli\watt\per\meter\squared} & Basal geothermal flux \\
$A$        & --        & Hemispherical bond albedo \\
$\varepsilon$ & --     & Thermal infrared emissivity \\
$\sigma$   & \SI{5.6704e-8}{\watt\per\meter\squared\per\kelvin\tothe{4}} & Stefan--Boltzmann constant \\
$\Rsolar$  & \SI{1361}{\watt\per\meter\squared} & Solar constant at \SI{1}{au} \\
$F_{\mathrm{sky}}$ & --  & Sky view factor (\cref{chap:phase2}) \\
$\phi$     & --        & Volumetric ice fraction (\cref{chap:phase3}) \\
\bottomrule
\end{tabularx}
\caption{Nomenclature used throughout the book.}
\label{tab:symbols}
\end{table}
